% Generated draft: requires mathematical and source review.
\documentclass[11pt]{article}
\usepackage[margin=25mm]{geometry}
\usepackage{fontspec}
\setmainfont{Latin Modern Roman}
\usepackage{amsmath,amssymb,mathtools,unicode-math}
\setmathfont{Latin Modern Math}
\usepackage{enumitem,needspace,xurl,hyperref}
\setlength{\parindent}{0pt}
\setlength{\parskip}{5pt}
\setlength{\emergencystretch}{3em}
\newcommand{\N}{\mathbb{N}}
\newcommand{\Z}{\mathbb{Z}}
\newcommand{\Q}{\mathbb{Q}}
\newcommand{\R}{\mathbb{R}}
\newcommand{\C}{\mathbb{C}}
\newcommand{\F}{\mathbb{F}}
\newcommand{\A}{\mathbb{A}}
\newcommand{\PP}{\mathbb{P}}
\newcommand{\E}{\mathbb{E}}
\newcommand{\eps}{\varepsilon}
\newcommand{\dd}{\,\mathrm{d}}
\begin{document}
\section*{UnsolvedMath ID: 30006593}
\textbf{Source problem code:} catalogue entry
\subsection{Definitions and mathematical statement}
Let \(\mathbb H=\{z\in\C:\operatorname{Im}z>0\}\) and \(q=e^{2\pi i z}\) for \(z\in\mathbb H\). Define the discriminant modular form and Ramanujan's tau function by
\[
\Delta(z)=q\prod_{m=1}^{\infty}(1-q^m)^{24}
=\sum_{n=1}^{\infty}\tau(n)q^n.
\]
Thus \(\tau(n)\), for each positive integer \(n\), is the coefficient of \(q^n\) in this expansion. Let \(\mathcal P=\{p\in\Z:p\ge 2\text{ and }p\text{ is prime}\}\).

The requested problem is to use double-Eisenstein-series representations of \(\Delta\) to establish
\[
\forall n\in\Z_{\ge 1},\qquad \tau(n)\ne 0,
\]
or, as the explicitly weaker alternative, to establish
\[
\#\{p\in\mathcal P:\tau(p)=0\}<\infty.
\]
The latter means that there exists an integer \(B\ge 1\) such that \(\tau(p)\ne 0\) for every prime \(p>B\); no bound is specified in the supplied record.

The original supplied question asks whether ``this formula'' can be applied to the non-vanishing problem of \(\tau(n)\). The record identifies double-Eisenstein-series representations as the intended method but does not supply the formula or define the relevant double Eisenstein series. Consequently, the non-vanishing targets are precise, while the particular representation and the intended argument remain unspecified. The supplied report assessment describes the proposed argument as heuristic; no proof is asserted here.
\subsection{Short English statement}
Use a double-Eisenstein-series representation of \(\Delta\) to prove that every positive-index Ramanujan coefficient \(\tau(n)\) is nonzero, or at least that only finitely many primes have \(\tau(p)=0\). The specific representation is absent from the supplied record, and its proposed application is described there as heuristic.
\subsection{Sources}
\begin{itemize}
\item Local UnsolvedMath registry: \texttt{data/unsolvedmath/problems.json}, with \texttt{catalog/problems.json} as the supplementary registry; permanent ID \(30006593\), exact source problem code \texttt{catalogue entry}. The supplied complete record is the basis of this draft; the local files were not independently inspected.
\item Supplied report citation: \emph{Multiple Zeta Values and Modular Forms} (2026), Oberwolfach Reports, report 7/2026, pp.~433--486, DOI \texttt{10.4171/owr/2026/7}. Source URL and supplied literature-source link, labeled ``Oberwolfach Report 7/2026'': \url{https://doi.org/10.4171/owr/2026/7}. The collection attributes the problem to workshop open-problem discussions; no individual proposer is supplied.
\item Supplied dated literature curation, checked \textbf{2026-08-22}: the record labels the problem open and states that the report presents this as an open application of a new formula to tau non-vanishing, with the proposed argument remaining heuristic. This is supplied dated curation, not an independently verified literature assessment. Its time-sensitive status and cited claims require independent verification.
\item The registry assigns the statement status \texttt{corrected\_verified} and says that the statement was reconstructed and checked against the official report. These are supplied provenance claims, not verification performed for this draft.
\end{itemize}
\end{document}
